% SOURCE_AUTHORITY: sga5_fr_workpass.tex, lines 8980--9010 (LF-normalized unit).
% SOURCE_SHA256: 791F4EFFC5E02832D5D77ED03518C8156D6F07E4C8238B03545DB93D883FBB28
\subsubsection*{3.2}
Sejam \(\ell\) um número primo \(\ne\car k\) e \(X\) um esquema liso de tipo finito sobre \(k\). Se \(Z\) é um ciclo algébrico de codimensão \(d\) de \(X\), na Exposição IV\footnote{Veja-se (SGA \(4^{1/2}\) Cycle).} foi associada a ele, para todo inteiro \(n\), uma classe de cohomologia
\[
\gamma_X^{(n)}(Z)\in H^{2d}\bigl(X,(\mu_X^{(n)})^{\otimes d}\bigr),
\]
onde \(\mu_X^{(n)}\) denota o \((\mathbb Z/\ell^{\,n+1}\mathbb Z)_X\)-módulo localmente livre das raízes \(\ell^{\,n+1}\)-ésimas da unidade. Por outro lado, foi observado que, se \(n'\ge n\), \(\gamma_X^{(n)}(Z)\) é a imagem de \(\gamma_X^{(n')}(Z)\) em
\[
H^{2d}\bigl(X,(\mu_X^{(n)})^{\otimes d}\bigr)
\]
pelo morfismo
\[
H^{2d}\bigl(X,(\mu_X^{(n')})^{\otimes d}\bigr)
\longrightarrow
H^{2d}\bigl(X,(\mu_X^{(n)})^{\otimes d}\bigr)
\]
deduzido da elevação à potência \(\ell^{\,n'-n}\):
\[
\mu_X^{(n')}\longrightarrow \mu_X^{(n)}.
\]
Associa-se assim a \(Z\) um elemento, denotado por \(\gamma_X(Z)\):
\[
\gamma_X(Z)\in H^{2d}(X,\mathbb Z_\ell(d)),
\]
onde \(\mathbb Z_\ell(d)\) se define em (1.3.4). Mais geralmente, se \(Z\) é um ciclo de \(X\), associa-se a ele por linearidade uma classe de cohomologia \(\ell\)-ádica em
\[
\bigoplus_d H^{2d}(X,\mathbb Z_\ell(d)),
\]
que é, de facto, o limite projetivo das classes de cohomologia associadas a \(Z\) nos
\[
\bigoplus_d H^{2d}\bigl(X,(\mu_X^{(n)})^{\otimes d}\bigr).
\]
